% Options for packages loaded elsewhere
% Options for packages loaded elsewhere
\PassOptionsToPackage{unicode}{hyperref}
\PassOptionsToPackage{hyphens}{url}
\PassOptionsToPackage{dvipsnames,svgnames,x11names}{xcolor}
%
\documentclass[
  12pt,
  openany]{book}
\usepackage{xcolor}
\usepackage[top=20mm,left=15mm,right=15mm,bottom=20mm]{geometry}
\usepackage{amsmath,amssymb}
\setcounter{secnumdepth}{5}
\usepackage{iftex}
\ifPDFTeX
  \usepackage[T1]{fontenc}
  \usepackage[utf8]{inputenc}
  \usepackage{textcomp} % provide euro and other symbols
\else % if luatex or xetex
  \usepackage{unicode-math} % this also loads fontspec
  \defaultfontfeatures{Scale=MatchLowercase}
  \defaultfontfeatures[\rmfamily]{Ligatures=TeX,Scale=1}
\fi
\usepackage{lmodern}
\ifPDFTeX\else
  % xetex/luatex font selection
\fi
% Use upquote if available, for straight quotes in verbatim environments
\IfFileExists{upquote.sty}{\usepackage{upquote}}{}
\IfFileExists{microtype.sty}{% use microtype if available
  \usepackage[]{microtype}
  \UseMicrotypeSet[protrusion]{basicmath} % disable protrusion for tt fonts
}{}
\makeatletter
\@ifundefined{KOMAClassName}{% if non-KOMA class
  \IfFileExists{parskip.sty}{%
    \usepackage{parskip}
  }{% else
    \setlength{\parindent}{0pt}
    \setlength{\parskip}{6pt plus 2pt minus 1pt}}
}{% if KOMA class
  \KOMAoptions{parskip=half}}
\makeatother
% Make \paragraph and \subparagraph free-standing
\makeatletter
\ifx\paragraph\undefined\else
  \let\oldparagraph\paragraph
  \renewcommand{\paragraph}{
    \@ifstar
      \xxxParagraphStar
      \xxxParagraphNoStar
  }
  \newcommand{\xxxParagraphStar}[1]{\oldparagraph*{#1}\mbox{}}
  \newcommand{\xxxParagraphNoStar}[1]{\oldparagraph{#1}\mbox{}}
\fi
\ifx\subparagraph\undefined\else
  \let\oldsubparagraph\subparagraph
  \renewcommand{\subparagraph}{
    \@ifstar
      \xxxSubParagraphStar
      \xxxSubParagraphNoStar
  }
  \newcommand{\xxxSubParagraphStar}[1]{\oldsubparagraph*{#1}\mbox{}}
  \newcommand{\xxxSubParagraphNoStar}[1]{\oldsubparagraph{#1}\mbox{}}
\fi
\makeatother


\usepackage{longtable,booktabs,array}
\newcounter{none} % for unnumbered tables
\usepackage{calc} % for calculating minipage widths
% Correct order of tables after \paragraph or \subparagraph
\usepackage{etoolbox}
\makeatletter
\patchcmd\longtable{\par}{\if@noskipsec\mbox{}\fi\par}{}{}
\makeatother
% Allow footnotes in longtable head/foot
\IfFileExists{footnotehyper.sty}{\usepackage{footnotehyper}}{\usepackage{footnote}}
\makesavenoteenv{longtable}
\usepackage{graphicx}
\makeatletter
\newsavebox\pandoc@box
\newcommand*\pandocbounded[1]{% scales image to fit in text height/width
  \sbox\pandoc@box{#1}%
  \Gscale@div\@tempa{\textheight}{\dimexpr\ht\pandoc@box+\dp\pandoc@box\relax}%
  \Gscale@div\@tempb{\linewidth}{\wd\pandoc@box}%
  \ifdim\@tempb\p@<\@tempa\p@\let\@tempa\@tempb\fi% select the smaller of both
  \ifdim\@tempa\p@<\p@\scalebox{\@tempa}{\usebox\pandoc@box}%
  \else\usebox{\pandoc@box}%
  \fi%
}
% Set default figure placement to htbp
\def\fps@figure{htbp}
\makeatother





\setlength{\emergencystretch}{3em} % prevent overfull lines

\providecommand{\tightlist}{%
  \setlength{\itemsep}{0pt}\setlength{\parskip}{0pt}}



 


\usepackage{fancyhdr}
\pagestyle{fancy}
\fancyhead[R]{ENSGSI 2AP - Mécanique des Solidés Déformables}

\fancyfoot[L]{Guillaume PRONOST | Fabio CRUZ}
\fancyfoot[C]{\thepage}
\fancyfoot[R]{ENSGSI - 2023/2024}
\renewcommand{\footrulewidth}{0.3pt}% default is 0pt
\setlength{\headheight}{15pt}
\usepackage{mathtools}
\usepackage[makeroom]{cancel}

\usepackage{fancyhdr} % en-têtes et pieds de page

\let\paragraph\oldparagraph
\let\subparagraph\oldsubparagraph

\usepackage[explicit]{titlesec} % titres de chapitres, sections, etc.




% Mise en forme des titres de chapitres : 

\titlespacing*{\chapter}{0pt}{-50pt}{5pt} 
\titleformat{\chapter}[display]
  {\LARGE\bfseries\sffamily}
  {}
  {12pt}
  {%
\begin{tcolorbox}[
	enhanced jigsaw,
	title={Mécanique du Point Matériel},
	fonttitle=\normalsize\sffamily,
	halign=center,
	top=\baselineskip,
	colback=red!10!white,
	colframe=red!50!black,
	boxrule=1.2pt,
	colbacktitle=red!50!white!75!black,
	attach boxed title to top center={yshift=-0.25mm-\tcboxedtitleheight/			2,yshifttext=2mm-\tcboxedtitleheight/2},
	boxed title style={enhanced,boxrule=0.5mm,
	frame code={ \path[tcb fill frame] ([xshift=-4mm]frame.west) -- 				(frame.north west)
    -- (frame.north east) -- ([xshift=4mm]frame.east)
    -- (frame.south east) -- (frame.south west) -- cycle; },
    interior code={ \path[tcb fill interior] ([xshift=-2mm]interior.west)
    -- (interior.north west) -- (interior.north east)
    -- ([xshift=2mm]interior.east) -- (interior.south east) -- (interior.south 		west)
    -- cycle;}  },
 	drop shadow]
    #1
\end{tcolorbox}
}

\renewcommand\thesection{\arabic{section}}
\makeatletter
\@ifpackageloaded{tcolorbox}{}{\usepackage[skins,breakable]{tcolorbox}}
\@ifpackageloaded{fontawesome5}{}{\usepackage{fontawesome5}}
\definecolor{quarto-callout-color}{HTML}{909090}
\definecolor{quarto-callout-note-color}{HTML}{0758E5}
\definecolor{quarto-callout-important-color}{HTML}{CC1914}
\definecolor{quarto-callout-warning-color}{HTML}{EB9113}
\definecolor{quarto-callout-tip-color}{HTML}{00A047}
\definecolor{quarto-callout-caution-color}{HTML}{FC5300}
\definecolor{quarto-callout-color-frame}{HTML}{acacac}
\definecolor{quarto-callout-note-color-frame}{HTML}{4582ec}
\definecolor{quarto-callout-important-color-frame}{HTML}{d9534f}
\definecolor{quarto-callout-warning-color-frame}{HTML}{f0ad4e}
\definecolor{quarto-callout-tip-color-frame}{HTML}{02b875}
\definecolor{quarto-callout-caution-color-frame}{HTML}{fd7e14}
\makeatother
\makeatletter
\@ifpackageloaded{caption}{}{\usepackage{caption}}
\AtBeginDocument{%
\ifdefined\contentsname
  \renewcommand*\contentsname{Table of contents}
\else
  \newcommand\contentsname{Table of contents}
\fi
\ifdefined\listfigurename
  \renewcommand*\listfigurename{List of Figures}
\else
  \newcommand\listfigurename{List of Figures}
\fi
\ifdefined\listtablename
  \renewcommand*\listtablename{List of Tables}
\else
  \newcommand\listtablename{List of Tables}
\fi
\ifdefined\figurename
  \renewcommand*\figurename{Figure}
\else
  \newcommand\figurename{Figure}
\fi
\ifdefined\tablename
  \renewcommand*\tablename{Table}
\else
  \newcommand\tablename{Table}
\fi
}
\@ifpackageloaded{float}{}{\usepackage{float}}
\floatstyle{ruled}
\@ifundefined{c@chapter}{\newfloat{codelisting}{h}{lop}}{\newfloat{codelisting}{h}{lop}[chapter]}
\floatname{codelisting}{Listing}
\newcommand*\listoflistings{\listof{codelisting}{List of Listings}}
\makeatother
\makeatletter
\makeatother
\makeatletter
\@ifpackageloaded{caption}{}{\usepackage{caption}}
\@ifpackageloaded{subcaption}{}{\usepackage{subcaption}}
\makeatother
\usepackage{bookmark}
\IfFileExists{xurl.sty}{\usepackage{xurl}}{} % add URL line breaks if available
\urlstyle{same}
% fallback for those not using the hyperref driver hyperxmp:
\makeatletter
\@ifundefined{xmpquote}{\newcommand{\xmpquote}[1]{#1}}{}
\makeatother
\hypersetup{
  pdftitle={Ethnographic Analysis of Urban Manufacturing Ecosystems},
  colorlinks=true,
  linkcolor={Maroon},
  filecolor={Maroon},
  citecolor={Blue},
  urlcolor={Blue},
  pdfcreator={LaTeX via pandoc}}


\title{Ethnographic Analysis of Urban Manufacturing Ecosystems}
\usepackage{etoolbox}
\makeatletter
\providecommand{\subtitle}[1]{% add subtitle to \maketitle
  \apptocmd{\@title}{\par {\large #1 \par}}{}{}
}
\makeatother
\subtitle{Internship proposition at the ERPI Laboratory}
\author{}
\date{Version 1 --- Updated on October 14, 2026}
\begin{document}
\frontmatter
\maketitle
\begin{abstract}
\textbf{Version 1}. \emph{Last updated on: Version 1 --- Updated on
October 14, 2026.}
\end{abstract}


\mainmatter
\chapter{Context of the Internship}\label{context-of-the-internship}

ERPI laboratory (\url{https://erpi.univ-lorraine.fr/}) is interested in
the development of methodologies, tools and case studies in the area of
innovation process management. One of the research axis concerns the
intersection on three major societal issues as follows:

\begin{enumerate}
\def\labelenumi{\arabic{enumi}.}
\tightlist
\item
  The circular economy of plastics as disruptor paradigm in political,
  industrial and societal perspective (\textbf{depascale2023?};
  \textbf{ramli2024?}),
\item
  The development of open design (\textbf{bakirlioglu2019?};
  \textbf{pacheco2025?}) including the open source technology as
  alternative approach of traditional technology development and,
\item
  The experimental development of the urban manufacturing approaches
  (\textbf{herrmann2020?}; \textbf{ijassi2025?}) (i.e.~Fablabs, Tiers
  Lieux, Recycleries and Matériauthèque) that usually are grassroot
  mouvements. These approaches are reframing the manufacturing networks
  at a local level.
\end{enumerate}

In this regard, ERPI Laboratory have been developing the concept of
``\emph{distributed recycling via additive manufacturing (DRAM)}''
(\textbf{cruzsanchez2020?}) as illustrated in Figure~\ref{fig-dram}. It
involves the technical feasibility of using recycled plastics into the
additive manufacturing printing process chain. Particularly, the use
systems based on material extrusion which are the most used in industry
and general public.

\begin{figure}

\centering{

\includegraphics[width=0.85\linewidth,height=\textheight,keepaspectratio,alt={alt}]{assets/DRAM.png}

}

\caption{\label{fig-dram}Distributed recycling via additive
manufacturing developed at the ERPI Laboratory}

\end{figure}%

The \textbf{Recovery (I)} phase addresses the logistical operations
required to collect plastic waste for reuse within DRAM systems. The
\textbf{Preparation (II)} phase encompasses the identification,
separation, sorting, size reduction, and cleaning of plastic waste to
ensure adequate quality. The \textbf{Compounding (III)} phase focuses on
the development of mono- and composite materials. After, the
\textbf{Feedstock (IV)} phase focuses on producing materials suitable
for printing, either in filament form for FFF or as particles for FGF.
The \textbf{Printing (V)} phase addresses application development and
process optimization for recycled materials. Finally, the
\textbf{Quality (VI)} phase involves multi-level technical
characterization of recycled materials and printed parts.

Therefore, three major scales are simultaneaosly studied, as illustrated
in Figure~\ref{fig-dram} :

\begin{itemize}
\tightlist
\item
  \textbf{Material level}: The study of strategies for post-consumer
  plastic recycling/upcycling to enhance the structural properties of
  plastic waste.
\item
  \textbf{Process level}: Open-source hardware (i.e additive
  manufacturing, manual injection, sheetpress etc) as a promising
  pathway to transform plastic waste into usable artifacts/products.
\item
  \textbf{Socio-technical level}: Urban manufacturing systems can
  provide a conducive environment for distributed upcycling, enabling
  the transformation of plastic waste into locally relevant products,
  closer to where the waste is generated.
\end{itemize}

For instance, using desktop-scale fused filament fabrication (FFF)
(\textbf{cruzsanchez2014?}) and fused granular fabrication (FGF)
(\textbf{alexandre2020?}), we have demonstrated their potential as
technical enablers for the mechanical recycling of mono-thermoplastic
waste, such as polylactic acid (PLA), from \emph{material}
(\textbf{cruz2015?}; \textbf{cruzsanchez2017?}), \emph{process}
(\textbf{cruzsanchez2014?}; \textbf{alexandre2020?}), and
\emph{socio-technical} perspectives (\textbf{santander2020?};
\textbf{caceres-mendoza2023?}). Recently, the work on multi-material
(\textbf{suescungonzalez2024?}; \textbf{gonzalez2025?}) have openend new
explorative paths that aims to see the \emph{upcycling}
(\textbf{riondet2022?}) potential of this approach.

\chapter{Main Issue}\label{main-issue}

This intership is placed at the \emph{Socio-technical level}. Urban
manufacturing is increasingly seen as a pathway to re-localize
production through cleaner and more sustainable processes
(\textbf{herrmann2019?}). It responds to major societal trends,
including urbanization, resource scarcity, circularity, and demand for
personalized production (\textbf{sajadieh2022?}). However, there is
still limited research on the types of products and systems that are
suitable for urban production environments. Integrating recycling and
manufacturing within cities has the potential to reduce transportation
emissions, support local economies, and foster citizen engagement in
sustainability practices.

Therefore, the major question to answer is :

\begin{itemize}
\tightlist
\item
  \emph{What socio-technical reconfigurations does the implementation of
  distributed recycling/upcycling processes bring about within a given
  urban territory, and under which social, material and institutional
  conditions can these processes take hold? (RQ3)}
\end{itemize}

Rather than treating social and environmental dimensions as side effects
of a technical system, the project approaches them as constitutive of
it: the collectives, skills, places and local policies through which
plastic waste comes to be qualified as a resource are part of what makes
distributed recycling possible, or not.

\section{Objective of the Internship}\label{objective-of-the-internship}

The internship takes place at an exploratory stage of the project, ahead
of the in-depth ethnographic fieldwork in Nancy that will be carried out
as part of a PhD (starting end of 2027). Its main goals are:

\begin{itemize}
\tightlist
\item
  To map the actors, networks and public action frameworks shaping
  citizen-led plastic upcycling in urban manufacturing settings, from
  the Nancy area (local waste governance, metropolitan authority, waste
  collection, reuse and repair organizations, public support for third
  places) to regional and national networks.
\item
  To build an analytical framework and a first set of contrasts that
  will help situate the Nancy case (Green Fablab, Octroi Nancy third
  place) in the subsequent fieldwork.
\end{itemize}

\chapter{Methodology}\label{methodology}

The internship will combine a literature review in political sociology,
the sociology of public action and the study of third places in France
(\textbf{gauthier2025?}), with a mapping of the actors and institutions
involved. The mapping will rely on document analysis (policy documents,
local waste plans, network resources) and on exploratory interviews with
regional and national networks (e.g., French Network of Fablabs RFFLabs,
Tiers-Lieux Grand Est, the Precious Plastic France network), as well as
with local institutional actors in the Nancy area. Initial contacts with
the Green Fablab and the Octroi Nancy third place will prepare the
in-depth fieldwork, without being the core of the internship. These
interviews are not intended to generalize, but to provide points of
contrast that will highlight what is specific to the Nancy
configuration.

\begin{tcolorbox}[enhanced jigsaw, arc=.35mm, bottomrule=.15mm, bottomtitle=1mm, breakable, colback=white, colbacktitle=quarto-callout-note-color!10!white, colframe=quarto-callout-note-color-frame, coltitle=black, left=2mm, leftrule=.75mm, opacityback=0, opacitybacktitle=0.6, rightrule=.15mm, title=\textcolor{quarto-callout-note-color}{\faInfo}\hspace{0.5em}{Note on scope}, titlerule=0mm, toprule=.15mm, toptitle=1mm]

This programme should be read as a horizon rather than a fixed set of
tasks; its actual scope will be adjusted with the intern according to
their profile and the progress of the work.

\end{tcolorbox}

The internship combines theoretical grounding and empirical
investigation.

The project will adopt a mixed qualitative and quantitative approach to
analyze urban upcycling systems, using as conceptual framework a Design
for Sustainability (DfS) perspective (\textbf{ceschin2019?}), focusing
on its socio-technical implications. First, a targeted literature review
in political sociology, public action, and the role of third places in
France (\textbf{gauthier2025?}) will establish an analytical framework
for understanding citizen-led initiatives. Assistant Professor Baptiste
Monsaingeon will support this task, bringing expertise in waste
sociology, particularly regarding citizen-led initiatives in France.

This will be complemented by in-depth ethnographic fieldwork in Nancy,
focusing on the Green Fablab and Octroi Nancy third place as a living
laboratory. Data collection will include participant observation,
semi-structured interviews, and participatory workshops. Ethical
considerations will be paramount, ensuring informed consent,
confidentiality, and respectful engagement with all participants. To
extend beyond the local scale, additional interviews will be conducted
with national networks (e.g., fablabs in Lorraine, Tiers Lieux Grand
Est, and Precious Plastic France network, French Network of Fablabs
RFFLabs), enabling comparative insights. The analysis will map
stakeholders' motivations, practices, constraints, and coordination
mechanisms across the ecosystem.

\section{Deliverable}\label{deliverable}

\begin{enumerate}
\def\labelenumi{\arabic{enumi}.}
\tightlist
\item
  A mapping report of the actors, networks and public action frameworks
  relevant to distributed plastic upcycling, from the Nancy area to the
  national level.
\item
  A firts scientific document presenting the analytical framework, the
  methodology and the first results, as a basis for the PhD fieldwork.
\end{enumerate}

\section{Profile}\label{profile}

\begin{itemize}
\tightlist
\item
  \textbf{Background}: Master M2 in Sociology of Innovation.
\item
  \textbf{Languages}: French and English.
\end{itemize}

\section{Modalities}\label{modalities}

\begin{itemize}
\tightlist
\item
  \textbf{Timming}: March 2027 -- September 2027
\item
  \textbf{Place}: ERPI, Université de Lorraine, Nancy, France.
\item
  \textbf{Salary}: Gratification for internships at the Université de
  Lorraine.
\end{itemize}

\section{Contacts}\label{contacts}

Send the CV and cover letter to:

\begin{itemize}
\tightlist
\item
  MCF Fabio A. Cruz Sanchez (fabio.cruz{[}at{]}univ-lorraine.fr)
\end{itemize}

\newpage

\chapter{References}\label{references}


\backmatter


\end{document}
